\section{Proof of Theorem~\ref{thm:main}(c)}\label{app:c}

Let $\lambda_d(u) = (d-1)\log u - (u - \kappa)^2/2$, so that $U$ has density $\propto e^{\lambda_d}$ on $(0, \infty)$. Throughout,
$\kappa = \sqrt{d/2} + \delta$, and every $O(\cdot)$ is uniform in $|\delta| \le B$ for any fixed $B$.

\begin{lemma}[The mode]\label{lem:mode}
$\lambda_d$ is strictly concave with unique maximiser $u^* = (\kappa + \sqrt{\kappa^2 + 4(d-1)})/2$. Let
$a_d = -\lambda_d''(u^*) = 1 + (d-1)/u^{*2}$. Then $1 \le a_d \le 2$ and
\begin{align*}
u^* &= \sqrt{2d} + \tfrac23\delta + \frac{4\delta^2 - 18}{27\sqrt{2d}} + O(1/d), &
2\kappa - u^* &= \tfrac43\delta + \frac{18 - 4\delta^2}{27\sqrt{2d}} + O(1/d),\\
a_d &= \tfrac32 - \frac{\sqrt2}{3}\,\frac{\delta}{\sqrt d} + O(1/d), &
\lambda_d'''(u^*) &= \frac{2(d-1)}{u^{*3}} = \frac{1}{\sqrt{2d}} + O(1/d).
\end{align*}
\end{lemma}

\begin{proof}
$\lambda_d'' = -1 - (d-1)/u^2 < 0$, and $u^{*2} = \kappa u^* + d - 1 \ge d - 1$ gives $a_d \le 2$. Next,
$\kappa^2 + 4(d-1) = \tfrac92 d + \sqrt{2d}\,\delta + \delta^2 - 4$; factoring out $\tfrac92 d$ and expanding the square
root to second order gives $\sqrt{\kappa^2 + 4(d-1)} = 3\sqrt{d/2} + \delta/3 + (8\delta^2 - 36)/(27\sqrt{2d}) + O(1/d)$,
hence $u^*$ and $2\kappa - u^*$. Then $u^{*2} = 2d + \tfrac{4\sqrt2}{3}\delta\sqrt d + O(1)$ gives $a_d$ and
$\lambda_d'''(u^*)$.
\end{proof}

\begin{lemma}[Laplace]\label{lem:laplace}
Let $Z_d = \sqrt{a_d}\,(U - u^*)$. Then $\sup_{|\delta| \le B}\sup_x|\Pr(Z_d \le x) - \Phi(x)| \to 0$.
\end{lemma}

\begin{proof}
Let $v = u - u^*$ and $\psi(v) = \lambda_d(u^* + v) - \lambda_d(u^*)$, so $\psi(0) = \psi'(0) = 0$ and $\psi''(0) = -a_d$. On
the window $|v| \le d^{1/8}$, $u \ge \sqrt{2d}/2$ for large $d$, so $|\lambda_d'''| \le 4\sqrt2/\sqrt d$ and
$\psi(v) = -a_dv^2/2 + O(d^{-1/8})$. Outside it, concavity gives $\psi(v) \le (|v|/h)\psi(\pm h)$ with $h = d^{1/8}$, and
$\psi(\pm h) \le -h^2/4$ since $a_d \ge 1$, so $\int_{|v| \ge h} e^\psi \le (8/h)e^{-h^2/4} \to 0$. Hence
$\int e^\psi = \sqrt{2\pi/a_d}\,(1 + o(1))$ and $\int_{-u^*}^{x/\sqrt{a_d}} e^\psi = \sqrt{2\pi/a_d}\,\Phi(x) + o(1)$
uniformly in $x$.
\end{proof}

\begin{lemma}[Next order]\label{lem:next}
There are $C$ and $d_0$ such that for $d \ge d_0$, $|\delta| \le B$ and $x \in \R$,
\[ \bigl|\Pr(Z_d \le x) - \Phi(x) + \varsigma_d(x^2 + 2)\varphi(x)\bigr| \le C/d, \qquad
   \varsigma_d = \frac{\lambda_d'''(u^*)}{6a_d^{3/2}}. \]
\end{lemma}

\begin{proof}
\emph{Window} $|v| \le \log d$: there $u \ge \sqrt d$ for $d \ge d_0$, so $|\lambda_d''''(u)| = 6(d-1)/u^4 \le 6/d$
and, with $b = \lambda_d'''(u^*) = O(d^{-1/2})$, Taylor's theorem gives $\psi(v) = -a_dv^2/2 + E(v)$ with
$E(v) = bv^3/6 + R_4(v)$ and $|R_4(v)| \le v^4/4d$. Since $|E| \le 1$ on the window for $d \ge d_0$,
$|e^E - 1 - E| \le E^2$, so $e^{\psi(v)} = e^{-a_dv^2/2}(1 + bv^3/6 + Q(v))$ with $|Q(v)| \le C(v^4 + v^6 + v^8)/d$.
\emph{Error}: since $a_d \ge 1$, $\int e^{-a_dv^2/2}|Q| = O(1/d)$, and extending the main term from the window to $\R$ costs
$O(e^{-(\log d)^2/4})$. \emph{Tails}: $\psi(\pm\log d) \le -(\log d)^2/4$ for $d \ge d_0$, so by concavity
$\int_{|v| \ge \log d} e^\psi \le (8/\log d)e^{-(\log d)^2/4}$, and $e^{-(\log d)^2/4} = d^{-(\log d)/4} = o(d^{-k})$ for every $k$.
\emph{Assemble}: $\int_{-u^*}^{y} e^\psi = \int_{-\infty}^{y} e^{-a_dv^2/2}(1 + bv^3/6)\dd v + O(1/d)$ uniformly in $y$; the
cubic term is odd, so the normaliser is $\sqrt{2\pi/a_d}\,(1 + O(1/d))$. Substitute $w = \sqrt{a_d}\,v$, use
$\int_{-\infty}^x w^3\varphi(w)\dd w = -(x^2 + 2)\varphi(x)$, and divide.
\end{proof}

\begin{proof}[Proof of Theorem~\ref{thm:main}(c)]
By Lemma~\ref{lem:mode}, $x_d(\delta) := (2\kappa - u^*)\sqrt{a_d} = \bar x_d(\delta) + O(1/d)$ and
$\varsigma_d = 1/(9\sqrt{3d}) + O(1/d)$, where $\bar x_d(\delta) = \sqrt{8/3}\,\delta + \sqrt3(3 - 2\delta^2)/(9\sqrt d)$;
the second term combines $\frac{18 - 4\delta^2}{27\sqrt2}$ with $-\tfrac43\delta\cdot\tfrac{\sqrt2}{9}\delta$. Since
$M_d = \Pr(Z_d \le x_d)$, Lemma~\ref{lem:next} and the Lipschitz continuity of $\Phi$ and of $(x^2 + 2)\varphi(x)$ give
\[ M_d(\sqrt{d/2} + \delta) = \Phi(\bar x_d(\delta)) - \frac{(\bar x_d^2 + 2)\varphi(\bar x_d)}{9\sqrt{3d}} + O(1/d) \]
uniformly in $|\delta| \le B$. Put $\delta_0 = z\sqrt{3/8}$, so $2\delta_0^2 = \tfrac34 z^2$, and
$\delta_\pm = \delta_0 + c/\sqrt d \pm K/d$. Expanding $\Phi$ about $z$, the coefficient of $1/\sqrt d$ in
$M_d(\sqrt{d/2} + \delta_\pm)$ is $\varphi(z)\bigl[\sqrt{8/3}\,c + (1 - z^2/4)/\sqrt3 - (z^2 + 2)/(9\sqrt3)\bigr]$, which
vanishes exactly at $c = c_\alpha = (13z^2 - 28)/(72\sqrt2)$. Then
$M_d(\sqrt{d/2} + \delta_\pm) - (1 - \alpha) = \pm K\sqrt{8/3}\,\varphi(z)/d + O(1/d)$, where the $O(1/d)$ does not
depend on $K$ for $d \ge \max(d_0, K^2)$. Choosing $K$ large and using Theorem~\ref{thm:main}(a),
$\sqrt{d/2} + \delta_- < \kappa_d < \sqrt{d/2} + \delta_+$.
\end{proof}

\section{Proofs for Section~\ref{sec:three}}\label{app:three}

\begin{proof}[Proof of Theorem~\ref{thm:far}]
Write $b_i = r_i - \varrho = -\vartheta_*\cdot a_i + O(1/\varrho)$. Then
\[ S = \sum_i r_i^2 - 2\varrho\,\mathcal D(x) + \mathcal Q(x), \qquad \mathcal Q = \sum_i w_i^2 - 2\sum_i b_iw_i , \]
and $\mathcal Q$ is Lipschitz on bounded sets and bounded in $C^2$ on compact sets avoiding the anchors.

(i) The Hessian of $\|x - a_i\|$ is $(I - u_iu_i^\top)/w_i$, positive semidefinite with kernel spanned by $u_i$, so
$\nabla^2\mathcal D(x)$ is positive definite unless the three $u_i$ are parallel, which happens only on a line through all
three anchors. On a segment at most two of the distances are affine, so $\mathcal D$ is strictly convex; it is coercive, so
$x_F$ is unique. The subdifferential of $\mathcal D$ at $a_i$ is $\sum_{j \ne i}u_{ij}$ plus the closed unit disc, so
$x_F = a_i$ exactly when $|\sum_{j \ne i}u_{ij}| \le 1$; as $|u_{ij} + u_{ik}| = 2\cos(\gamma_i/2)$ for the angle $\gamma_i$ at
$a_i$, this is $\gamma_i \ge 120^\circ$. If every angle is below $120^\circ$, then $\nabla\mathcal D(x_F) = 0$ and
$\nabla^2\mathcal D(x_F) \succ 0$. On a closed disc $B$ about $x_F$ avoiding the anchors, $\nabla S/2\varrho = -\nabla\mathcal D +
O(1/\varrho)$ in $C^1$, so for large $\varrho$ $\nabla S$ has exactly one zero $x_\varrho$ in $B$, with $x_\varrho - x_F = O(1/\varrho)$, and
$\nabla^2S(x_\varrho)/2\varrho = -\nabla^2\mathcal D(x_F) + o(1)$ is negative definite. As $\nabla\mathcal D(x_F) = 0$,
$\mathcal D(x_\varrho) = \mathcal D(x_F) + O(\varrho^{-2})$, so $S(x_\varrho) - S(a_i) = 2\varrho(\mathcal D(a_i) - \mathcal D(x_F)) + O(1) > 0$.
At the anchors, $r_i/\varrho \to 1$ and $v_i/\varrho \to -\sum_{j \ne i}u_{ij}$, so the test of Remark~\ref{rem:three} becomes
$1 > |\sum_{j \ne i}u_{ij}|$ in the limit: it fails strictly where the angle is below $120^\circ$ and holds strictly where it
exceeds $120^\circ$. Finally, on a fixed disc around the anchors, $\nabla S/2\varrho = -\nabla\mathcal D + O(1/\varrho)$ away from
the anchors, $\nabla\mathcal D$ vanishes only at $x_F$, and near $a_i$, $\nabla\mathcal D = u_i + \sum_{j \ne i}u_j$ with
$\bigl||\sum_{j \ne i}u_j| - 1\bigr|$ bounded below, so $|\nabla\mathcal D| \ge c_1 > 0$ there while
$|\nabla\mathcal Q|/2\varrho = O(1/\varrho)$. So $x_\varrho$ is the only smooth critical point near the anchors when every angle
is below $120^\circ$, and there is none when one exceeds it.

(ii) Write $x = (\varrho + \tau)\vartheta$ with $|\vartheta| = 1$, and fix a small $c_2 > 0$. For $M_0 \le |x| \le (1 - c_2)\varrho$,
with $M_0$ fixed, all residuals $e_i$ are negative, and for $|x| \ge (1 + c_2)\varrho$ all are positive; since
$\partial_\tau w_i > 0$ there, $\partial_\tau S = 2\sum_i e_i\partial_\tau w_i \ne 0$. In the annulus
$(1 - c_2)\varrho \le |x| \le (1 + c_2)\varrho$, $w_i = \varrho + \tau - \vartheta\cdot a_i + O(1/\varrho)$ and
$r_i = \varrho - \vartheta_*\cdot a_i + O(1/\varrho)$, so $e_i = \tau + (\vartheta_* - \vartheta)\cdot a_i + O(1/\varrho)$, and
$\partial_\tau^2S \ge 5$ for small $c_2$ and large $\varrho$. So $\partial_\tau S = 0$ defines $\tau^*(\vartheta)$ uniquely, and
the critical points of $S$ there are those of $\vartheta \mapsto S((\varrho + \tau^*(\vartheta))\vartheta)$, which converges in
$C^2$ of the circle to $\min_\tau\sum_i(\tau + (\vartheta_* - \vartheta)\cdot a_i)^2 = G(\vartheta)$. Non-degenerate critical
points persist one-to-one, with the radial direction a minimum. On the circle,
$G(\vartheta) = \vartheta^\top M\vartheta - 2(M\vartheta_*)^\top\vartheta + \mathrm{const}$, whose critical points solve
$(M - \mu I)\vartheta = M\vartheta_*$, $|\vartheta| = 1$; there are four exactly when $M\vartheta_*$ lies inside the astroid
$|c_1'|^{2/3} + |c_2'|^{2/3} = (\lambda_1 - \lambda_2)^{2/3}$, the evolute of the corresponding ellipse, in coordinates
$(c_1', c_2')$ along the eigenvectors, and two otherwise. The left side of the stated inequality is smallest at
$\vartheta_*$ along the second eigenvector, where it is $\lambda_2^{2/3}$, so four needs $\lambda_1 > 2\lambda_2$.

(iii) By (i) and (ii) the critical points are $x^*$, one saddle and $x_\varrho$. At an anchor where the test fails, the
directional derivative $-2(r_i - v_i\cdot\vartheta)$ is negative on an open arc of directions and positive on the
complementary arc, so a vector field decreasing $S$ can be built near it as near a regular point, and the sublevel sets
change topology only at the three critical depths: an arc around $x^*$, a ring once the two ends meet at the saddle,
whose depth tends to $\frac12\max G$ by (ii), and a disc once the ring closes at $x_\varrho$. Finally
$S(x_F) = \sum_i(r_i - \|x_F - a_i\|)^2$, $\nabla S(x_F) = O(1)$ and $x_\varrho - x_F = O(1/\varrho)$ give
$S(x_\varrho) = S(x_F) + O(1/\varrho)$.
\end{proof}

\begin{proof}[Proof of Proposition~\ref{prop:birth}]
Here $r_1 = \varrho + 1$, $r_2 = r_3 = (\varrho^2 - \varrho + 1)^{1/2}$, $\|a_1 - a_j\| = \sqrt3$ and $u_{12} + u_{13} = (0, \sqrt3)$,
so $v_1 = (\sqrt3 - r_2)(0, \sqrt3)$ and, for $\varrho \ge 2$, where $r_2 \ge \sqrt3$, $|v_1| = \sqrt3\,r_2 - 3$. The test
$r_1 > |v_1|$ reads $(\varrho + 4)^2 > 3(\varrho^2 - \varrho + 1)$, that is, $(2\varrho - 13)(\varrho + 1) < 0$. For the other two
anchors the test $r_2 > |v_2|$ changes sign once on $[2, \infty)$, at $\varrho = 2.7227$ (computed).

On the axis put $x = (0, t)$, $-\frac12 < t < 1$. Then $w_1 = 1 - t$ and $w_2 = w_3 = w(t) = (t^2 + t + 1)^{1/2}$, with
$w' = (2t + 1)/2w$, $w'' = \frac34 w^{-3} > 0$ and $w''' = -\frac98(2t + 1)w^{-5} < 0$, and $f(t) = S(0, t)$ has
$f'(t)/2 = k(t) = 3t + \varrho + 1 - 2r_2w'(t)$. So $k'' = -2r_2w''' > 0$: $k$ is convex, with $k(-\frac12) = \varrho - \frac12 > 0$
and $k(1) = \varrho + 4 - \sqrt3\,r_2$, which is positive for $\varrho < 13/2$ and negative beyond. A double zero needs
$k = k' = 0$. Now $k' = 0$ is $w^3 = r_2/2$; with $y = 2t + 1$ and $Y = 2\varrho - 1$ it turns $k = 0$ into $y^3 = Y$, and since
$w^2 = (y^2 + 3)/4$ and $r_2^2 = (Y^2 + 3)/4$, with $q = Y^{2/3}$ this gives $(q + 3)^3 = 4(q^3 + 3)$, that is,
$(q + 1)^2(q - 5) = 0$. So the only double zero is at $\varrho_0 = (1 + 5\sqrt5)/2$, $t_0 = (\sqrt5 - 1)/2$, where $w = \sqrt2$
and $r_2 = 4\sqrt2$; it is a non-degenerate fold, since $k''(t_0) > 0$ and $\partial_\varrho k(t_0) = -\frac9{16}$. The number of
zeros of $k$ can change only at a double zero or when a zero crosses an end of the interval, and $k > 0$ on the whole
interval at $\varrho = 2$; so $k$ has no zero for $2 \le \varrho < \varrho_0$, two zeros $t_1 < t_2$ for $\varrho_0 < \varrho < 13/2$, and,
since $k(1) < 0$ beyond, exactly one for $\varrho > 13/2$, where $t_2$ has left through $t = 1$, the anchor $a_1$. At $t_1$,
$f$ has a local maximum, and at $t_2$ a local minimum. Across the axis,
\[ \partial_x^2S(0, t) = 2\Bigl[1 - \frac{r_1}{w_1} + 2\Bigl(\frac{3}{4w^2} + (w - r_2)\frac{(t + \frac12)^2}{w^3}\Bigr)\Bigr]
   < 2\Bigl[3 - \frac{\varrho + 1}{1 - t}\Bigr] < 0 \]
for $\varrho \ge 7/2$, since $w \le \sqrt3 \le r_2$ and $1 - t < \frac32$. So $t_1$ is a strict local maximum of $S$ and $t_2$ a
saddle. As $\varrho \to \infty$, $k/\varrho \to 1 - 2w'$, whose zero on the interval is $t = 0$, the centre.
\end{proof}

\section{A two-variable Laplace lemma}\label{app:laplace2}

\begin{lemma}\label{lem:laplace2}
Let $\mathcal K$ be a compact set of pairs (non-degenerate scaled triangle, $\sigma$). Under the posterior of
Lemma~\ref{lem:conc}, $\xi = \sqrt d\,(\tilde e - \tilde e^*)$ converges in total variation to $N(0, H^{-1})$, uniformly on
$\mathcal K$.
\end{lemma}

\begin{proof}
By the implicit function theorem $\tilde e^*$ and $H$ are continuous on $\mathcal K$, so there are $\nu, C, \bar\lambda > 0$ such
that, uniformly: the $\nu$-ball about $\tilde e^*$ lies in $\tilde F$; on it $|\nabla^3\mathcal E_\sigma| \le C$ and
$\psi = \tilde w_1\tilde w_2\tilde h^{-3}$ satisfies $1/C \le \psi \le C$ and $|\nabla\psi| \le C$; and
$I/\sigma_{\max}^2 \preceq H \preceq \bar\lambda I$. Let $\varrho_d = d^{-3/8}$.

\emph{Window} $|\tilde e - \tilde e^*| \le \varrho_d$, i.e.\ $|\xi| \le d^{1/8}$. By Taylor's theorem
$d(\mathcal E_\sigma(\tilde e) - \mathcal E_\sigma(\tilde e^*)) = -\xi^\top H\xi/2 + O(d\varrho_d^3)$ and $\psi(\tilde e) = \psi(\tilde e^*)(1 + O(\varrho_d))$,
so the density of $\xi$ is $\propto e^{-\xi^\top H\xi/2}(1 + o(1))$ there, uniformly.

\emph{Tails.} On the sphere $|\tilde e - \tilde e^*| = \varrho_d$, $\mathcal E_\sigma - \mathcal E_\sigma^* \le -\varrho_d^2/2\sigma_{\max}^2 + C\varrho_d^3$,
so by concavity along rays from $\tilde e^*$, $\mathcal E_\sigma(\tilde e) - \mathcal E_\sigma^* \le -c\varrho_d|\tilde e - \tilde e^*|$ outside the
window, with $c = 1/4\sigma_{\max}^2$, for large $d$. To avoid the singular factor $\tilde h^{-3}$, write the integrand of
Lemma~\ref{lem:conc} as $e^{(d-3)\mathcal E_\sigma}\,\tilde w_1\tilde w_2\,e^{-3|\tilde e|^2/2\sigma^2}$ and apply the bound to
$(d-3)\mathcal E_\sigma$. The tail mass relative to $e^{d\mathcal E_\sigma^*}$ is then at most
$e^{-3\mathcal E_\sigma^*}e^{-c(d-3)\varrho_d^2}\int\tilde w_1\tilde w_2 e^{-3|\tilde e|^2/2\sigma^2}\dd\tilde e = O(e^{-c'd^{1/4}})$, while the
window carries $\psi(\tilde e^*)(2\pi/d)\det(H)^{-1/2}(1 + o(1))$. So the tail is negligible, and the Gaussian mass of
$|\xi| > d^{1/8}$ tends to $0$ as well.
\end{proof}

\section{Proofs for Section~\ref{sec:high}}\label{app:high}

\begin{proof}[Proof of Lemma~\ref{lem:conc}]
Substituting $e = \sqrt d\,\tilde e$ in Lemma~\ref{lem:residual}(ii), all powers of $\sqrt d$ are constants, and
\[ e^{-|e|^2/2\sigma^2}w_1w_2h^{d-3} \propto e^{-d|\tilde e|^2/2\sigma^2}\,\tilde w_1\tilde w_2\,\tilde h^{d-3}, \]
which is the stated form. By Lemma~\ref{lem:weight}(i), $\tilde h^2$ is a constant times $\tilde s'\tilde\ell_*\tilde\ell_1\tilde\ell_2$, four
functions that are affine in $\tilde e$ and positive on $\tilde F$; so $\log\tilde h$ is concave, $\mathcal E_\sigma$ is strictly
concave, and $\mathcal E_\sigma \to -\infty$ where some $\tilde\ell_k \to 0$ and, since $\log\tilde h = O(\log|\tilde e|)$, at
infinity. Its superlevel sets are therefore compact subsets of $\tilde F$, and the maximiser exists, is interior and is
unique. At $0$ the quadratic term has zero gradient, and the component of $\nabla\log\tilde h(0)$ along
$(1,1)/\sqrt2$ is $\tfrac{\sqrt2}{4}(1/\tilde s + 1/\tilde t_*) > 0$, so $\tilde e^* \ne 0$. For a symmetric triangle the swap
$\tilde e_1 \leftrightarrow \tilde e_2$ preserves $\tilde F$ and $\mathcal E_\sigma$, so it fixes the unique maximiser, and $u > 0$
because $\mathcal E_\sigma$ is concave along the diagonal with positive slope at $0$.
\end{proof}

\begin{proof}[Proof of Lemma~\ref{lem:cdepth}]
Since $\log\tilde h$ is concave, $H \succeq I/\sigma^2$. Differentiating $\nabla\log\tilde h(\tilde e^*) = \tilde e^*/\sigma^2$ in
$\sigma$ (implicit function theorem) gives the derivative of $\tilde e^*$, and $\delta^{*\prime} = \tilde e^*\cdot\dd\tilde e^*/\dd\sigma$.
For $\sigma < \sigma'$ with maximisers $e, e'$, adding $\mathcal E_\sigma(e) \ge \mathcal E_\sigma(e')$ and
$\mathcal E_{\sigma'}(e') \ge \mathcal E_{\sigma'}(e)$ gives $(|e'|^2 - |e|^2)(\sigma^{-2} - \sigma'^{-2}) \ge 0$; equality would make
$e'$ maximise $\mathcal E_\sigma$ as well, so $e' = e$ and $e/\sigma^2 = e/\sigma'^2$, contradicting $e \ne 0$. As $\sigma \to 0$,
$\mathcal E_\sigma(\tilde e^*) \ge \mathcal E_\sigma(0)$ gives $|\tilde e^*|^2 \le 2\sigma^2(\log\tilde h(\tilde e^*) - \log\tilde h(0)) \to 0$.
If $\tilde e^*$ stayed bounded as $\sigma \to \infty$, then $\nabla\log\tilde h(\tilde e^*) = \tilde e^*/\sigma^2 \to 0$; but the
derivative of $\log\tilde h$ along $(1,1)/\sqrt2$ is $\tfrac12\bigl(1/(\sqrt2\,\tilde s') + 1/\tilde\ell_*\bigr)$, which is bounded
below on bounded sets.
\end{proof}

\begin{proof}[Proof of Theorem~\ref{thm:high}]
Let $\xi = \sqrt d(\tilde e - \tilde e^*)$. By Lemma~\ref{lem:laplace2}, $\xi$ converges to $N(0, H^{-1})$ in total
variation, uniformly. Since $\Delta = d|\tilde e|^2/2$, $(\Delta - d\delta^*)/\sqrt d = \tilde e^*\cdot\xi + |\xi|^2/2\sqrt d$. The first
term tends to $N(0, g^2)$; on $|\xi| \le M$ the event $\{\Delta \le \delta\}$ lies between two half-planes in $\xi$ shifted by
$\pm M^2/2\sqrt d$, and letting $M \to \infty$ slowly gives
\[ \sup_t\bigl|\Pi\bigl((\Delta - d\delta^*)/\sqrt d \le t\bigr) - \Phi(t/g)\bigr| \to 0 . \]
Since $g \ge |\tilde e^*|\,\lambda_{\max}(H)^{-1/2}$ is bounded
below, the quantiles converge, and the HPD region $\{\Delta \le \Delta_\alpha\}$, a one-sided event of mass $1 - \alpha$, gives
$\Delta_\alpha = d\delta^* + z\sqrt d\,g + o(\sqrt d)$. A transition of scaled depth $\tilde\Delta_c$ has depth $d\tilde\Delta_c$ and
occurs when $\Delta_\alpha(\sigma) = d\tilde\Delta_c$ (Theorem~\ref{thm:depth} and Proposition~\ref{prop:axis}), i.e.\ when
$\delta^*(\sigma) = \tilde\Delta_c - zg/\sqrt d + o(d^{-1/2})$; invert with Lemma~\ref{lem:cdepth}.
\end{proof}

\begin{proof}[Proof of Corollary~\ref{cor:closed}]
By Lemma~\ref{lem:conc}, $\tilde e^* = (u, u)$, where $\tilde w = 1 + u$, $\tilde h^2 = \tilde w^2 - \zeta^2$ and
$\mathcal E_\sigma = \tfrac12\log(\tilde w^2 - \zeta^2) - u^2/\sigma^2$, so stationarity gives $\sigma^2 = 2u\tilde h^2/\tilde w$ and
$\delta^* = u^2$. The tangent lengths are $\tilde t_* = 1 - \zeta$ and $\tilde t_1 = \tilde t_2 = \zeta$, so the inner contact comes first
exactly when $\theta > 60^\circ$. For the inner contact $u = 1 - \zeta$, $\tilde w = 2 - \zeta$ and $\tilde h^2 = 4(1 - \zeta)$; for the
outer contacts $u = \zeta$, $\tilde w = 1 + \zeta$, $\tilde h^2 = 1 + 2\zeta$. By symmetry the diagonal is an eigenvector of $H$, with
eigenvalue the curvature of $-\mathcal E_\sigma$ along arc length $\sqrt2\,u$,
$\lambda_\parallel = \tfrac12(\tilde w^2 + \zeta^2)/(\tilde w^2 - \zeta^2)^2 + 1/\sigma^2 = (\zeta^2 - 4\zeta + 6)/(16(1-\zeta)^2)$ at the inner contact;
so $g^2 = 2u^2/\lambda_\parallel = 32(1-\zeta)^4/(\zeta^2 - 4\zeta + 6)$ and Theorem~\ref{thm:high} gives the coefficient; at the outer
contacts the same formula gives $\lambda_\parallel = (2\zeta^3 + 4\zeta^2 + 4\zeta + 1)/(2\zeta(1+2\zeta)^2)$ and
$g^2 = 4\zeta^3(1+2\zeta)^2/(2\zeta^3 + 4\zeta^2 + 4\zeta + 1)$. For one range,
$\mathcal E_\sigma(\tilde e) = \log(1 + \tilde e) - \tilde e^2/2\sigma^2$ in one variable, the hole fills at depth $r^2/2$, i.e.\
$\tilde e^* = 1$, so $\sigma^2 = \tilde e^*(1 + \tilde e^*) = 2$, $H = \tfrac14 + \tfrac12$ and $g^2 = 4/3$; and
$\sqrt d/(\sqrt{d/2} + z\sqrt{3/8}) = \sqrt2 - 2z\sqrt{3/8}/\sqrt d + O(1/d)$.
\end{proof}

\section{Computation and validation}\label{app:comp}

\paragraph{One range} Substituting $v = u - \kappa$ and expanding $(v + \kappa)^{d-1}$ binomially writes both integrals
in $M_d$ as finite sums of truncated normal moments, available in closed form through $\Phi$, $\varphi$ and a recursion.
These sums alternate in sign and cancel heavily for large $d$, so they are evaluated in 60-digit arithmetic
\citep{mpmath}.
$\kappa_d(\alpha)$ is found by bisection on $M_d - (1 - \alpha)$, which is monotone by Theorem~\ref{thm:main}(a), to a
bracket shorter than $10^{-15}$. Every entry of Table~\ref{tab:kappa} was cross-checked by adaptive quadrature of the
defining integrals at 30 digits: at the computed $\kappa_d$ the quadrature value of $|M_d - (1 - \alpha)|$ is below
$10^{-16}$ for all 20 entries. The random-anchor probability is
$P_d(\sigma) = \int_0^\infty f_{\chi_d}(t)\,\Phi\bigl(t/\sigma - \kappa_d\bigr)\dd t$, by adaptive quadrature.

\paragraph{Two ranges: exact thresholds} Every two-range noise level in the paper is computed by quadrature, with no
grid. By Lemma~\ref{lem:residual}, $\Pi_\sigma(\Delta \le \delta) = N(\sqrt{2\delta})/N(\infty)$ with
$N(R) = \int_0^R e^{-u^2/2\sigma^2}u\,\mathcal J(u)\dd u$, where $\mathcal J(u)$ is the integral of $J$ over the arcs of the circle
$|e| = u$ that lie in $F$; the disc boundary is exact in these coordinates. By Lemma~\ref{lem:weight}(i), $J$ is a
product of powers of $w_1, w_2, s'$ and the three wall distances. We evaluate the wall distances in a cancellation-free
product form and remove every endpoint power-law factor by the substitution $t = t_{\mathrm{end}} \pm cv^2$, in angle and in
radius, so all integrands are smooth; composite Gauss--Legendre rules then converge fast. Doubling the nodes changes
every computed threshold by less than $10^{-9}$ (relative) for $d \ge 3$ and $5\cdot10^{-6}$ for $d = 2$, where $J$ has
inverse-square-root singularities at the walls. A transition of depth $\Delta_c$ occurs at the root of
$\Pi_\sigma(\Delta \le \Delta_c) = 1 - \alpha$, unique by Theorem~\ref{thm:depth}, found by Brent's method. This gives the
ladder of Section~\ref{sec:two} and the transition curves of Figure~\ref{fig:phase}, the exact column of
Table~\ref{tab:sharp} and the masses outside the disc at $\sigma_1$, and the finite-$d$ values of Figure~\ref{fig:high}. The
limits and coefficients of Corollary~\ref{cor:closed} agree with a direct two-dimensional maximisation of $\mathcal E_\sigma$
and a finite-difference Hessian to five decimals.

\paragraph{Two ranges: independent check of the topology} The topology is proved (Proposition~\ref{prop:axis}); the
computation below only checks it against brute force, under a criterion fixed in advance. For every crossing angle
$15^\circ, 20^\circ, \dots, 165^\circ$ and $d \in \{2, 3, 4, 6, 10\}$ we take the exact transition levels, and sample one noise
level below the first, one above the last, and the geometric midpoint of every interval between consecutive transitions
whose endpoints differ by at least $4\%$ in $\sigma$. At each such level we compare the Betti numbers predicted by
Proposition~\ref{prop:axis} with (a) the Betti numbers of the grid region in the half-plane ($900 \times 450$ cells, HPD
level from the sorted cell masses as in \citet{hyndman1996}), computed as a cubical complex and lifted by
Lemma~\ref{lem:lift}: $536$ of $536$ agree; and (b) for $d = 3$, at six angles and for the asymmetric example, the Betti
numbers of the full three-dimensional region on a $170^3$ grid, computed by GUDHI \citep{gudhi2014} without using
Lemma~\ref{lem:lift}: $28$ of $28$ agree. The $79$ intervals narrower than $4\%$ are all two-cavity phases, at crossing
angles of at most $25^\circ$ or at least $105^\circ$; the five widest of them, sampled on a $2400 \times 1200$ grid, also
agree. Proposition~\ref{prop:axis} was further checked on $12$ random triangles in $d = 2$ and $3$ at nine noise levels
each, discarding levels within $2\%$ of a critical value: all $210$ remaining cases agree. The three-dimensional figures
are isosurfaces of the posterior at the HPD level, extracted by marching cubes \citep{lorensen1987} from $170^3$--$180^3$
grids, and each panel's Betti numbers were confirmed by GUDHI.

\paragraph{Unequal noise} Theorem~\ref{thm:unequal} was checked against brute-force Betti numbers in the plane on adaptive
grids ($900$ to $1100$ cells per axis) for five fixed geometries and ten random geometries with random weights, at noise
levels whose HPD depth is more than $3\%$ from every critical depth: all agree. Its closed forms and the classification in (iii) were checked on
$20\,000$ random cases (noise-free point uniform on $[-3, 3] \times [0.05, 3]$, $\log \omega_1, \log \omega_2$ uniform on
$[-1.5, 1.5]$): the depths agree with direct minimisation of $\Delta$ over each edge, annihilation occurred in $65.5\%$ of
cases, never twice, always at the noisier anchor and never first. Corollary~\ref{cor:anchortest} and the two-cavity
criterion of Remark~\ref{rem:unequal} were checked on $4000$ further random cases, with no exceptions.

\paragraph{Up to $d$ ranges and the unequal-noise certificate} Theorem~\ref{thm:nranges} was checked for $n = d = 3$,
with anchors $(1, 0, 0)$, $(-\frac12, \pm\frac{\sqrt3}{2}, 0)$ and $x^* = (0.2, 0.1, 0.5)$: $\Delta_{\mathcal A} = 0.021605$,
reached by $\Delta_\alpha$ at $\sigma \approx 0.0726$, and counting components on a $281^3$ grid gives two pieces at
$0.5$, $0.8$ and $0.95$ times that level and one piece at $1.05$, $1.2$ and $1.6$ times it. For
Theorem~\ref{thm:staircaseW}, both certificates were evaluated in $36$ cases (four geometries, three noise ratios,
$d = 2, 3, 4$; $\beta = 1/\sqrt d$), and the posterior mass outside $|\bar e| < R_W$ was computed at $\sigma_1$ on a
$1600 \times 800$ half-plane grid with weight $\rho^{d-2}$: it lies between $9.4\cdot10^{-4}$ and $1.0\cdot10^{-2}$,
against the allowed $0.05$. The true first transition, the root of that mass minus $\alpha$, gives the ratios quoted
in \S\ref{sec:staircase}. The $d = 2$ certificate also holds at $\alpha = 0.1$, $0.3$ and $0.5$ in the cases tried.

\paragraph{Coverage} For Table~\ref{tab:coverage}, each data set's posterior is evaluated on a $600 \times 600$ grid over the
square of half-width $\max_i r_i + 1.2 + 5\sigma$ centred between the anchors, and the HPD region is formed from the
sorted cell masses. The mode is found by the Nelder--Mead method started from the grid maximum or its mirror image, and
the Hessian of $S$ there by central differences; it was positive definite in all $12\,000$ data sets. The mixture
places Gaussians with this covariance at the mode and its mirror image in the anchor line, with weights $\frac12$
(one Gaussian if the mode is on the line), and its HPD region is formed from sorted cell masses on the same grid. The moment-matched
Gaussian uses the grid mean and covariance. Each noise level uses its own fixed random seed, and rerunning reproduces
every entry exactly.
